\section{Introduction}

Continuous near-surface meteorological observations are central to Antarctic climate analysis, model evaluation, and field operations. Automatic weather stations (AWSs) extend coverage beyond staffed stations, but the same remoteness that makes them valuable also makes their records incomplete. The AntAWS compilation aggregates Antarctic AWS observations at 3-hourly, daily, and monthly resolutions and documents the role of AWS networks in resolving remote Antarctic conditions \citep{wang2023antaws}. In the selected 3-hourly stations used in this work, missingness is not well described by independent random points: wind speed and relative humidity contain long contiguous gaps, and several stations exhibit outage blocks lasting days to weeks.

The task is imputation, not forecasting: the model reconstructs missing station observations within a window while leaving available observations unchanged. ERA5 reanalysis provides a physically coherent meteorological background \citep{hersbach2020era5}, but using it naively can be counterproductive: when AWS observations are present, ERA5 often duplicates already available temperature and pressure information; when relative humidity is biased, always-on fusion can inject error. We therefore use ERA5 as a conditional prior that is activated by missingness.

For reproducibility, the study uses ECBIT, the ERA5-conditioned block imputation transformer, as a lightweight dual-stream architecture for sparse Antarctic AWS time series. One stream encodes normalized AWS values, observation masks, and time features; the other encodes aligned ERA5 variables. A conditioning module projects ERA5 tokens into the AWS embedding space, and a missing-variable gate restricts ERA5 influence to channels that contain missing positions. Loss is computed only on artificially hidden positions that were genuinely observed in the raw record, so the benchmark and label support remain fixed while ERA5 access, fusion form, and mask curriculum are varied independently.

The current benchmark contains 32 selected AntAWS stations with at least 10 years of record and sufficient multi-variable completeness. Five geographically diverse stations are reserved for held-out generalization. ERA5 point time series are downloaded and aligned to the exact AntAWS 3-hourly timestamps, producing 41,388 sparse 168-step windows. The experimental matrix covers interpolation, last-observation-carried-forward, ERA5 direct substitution, SAITS \citep{du2023saits}, an iTransformer-style imputation baseline \citep{liu2024itransformer}, ERA5-augmented SAITS/iTransformer controls, and ECBIT ablations.

The contributions are:
\begin{itemize}
  \item an AntAWS-derived benchmark of 32 selected 3-hourly stations with aligned ERA5 point series, sparse observation masks, realistic block-missing simulation, and held-out station splits;
  \item quantitative evidence that ERA5 conditioning reduces block-trained MAE from 0.346 to 0.258, with a relative benefit that grows from 13.8\% at 24 h to 26.4\% at 216 h, concentrated in temperature, wind speed, and specific humidity;
  \item a failure-mode curriculum study showing that MCAR-trained ERA5 models transfer poorly to block tests, reaching 0.407 MAE, whereas block-trained ERA5 models reach 0.258;
  \item fair ERA5-augmented baseline tests showing that augmented SAITS, augmented iTransformer, gated ECBIT, and concat ECBIT are statistically indistinguishable in aggregate MAE, placing the empirical emphasis on failure-mode matching and ERA5 conditioning rather than fusion architecture.
\end{itemize}
